% Tableau des labs du cours IIN dans NDAP SUKULU (propre au cours).
Sous le code \codeUE, NDAP SUKULU propose quatorze fiches d'auto-évaluation,
une par chapitre, prologue compris, et dix laboratoires. Tous portent sur
l'affaire MBONGO, le cas fil rouge du cours. Ils se font avec un terminal
simple et un tableur, sans outil spécialisé, et durent de quarante-cinq à
soixante minutes. Les laboratoires d'analyse sur pièces figurent dans le
cours de suite, Investigation Numérique Avancée (UE~SEC4052).

\begin{tableaucours}{lp{8.6cm}p{1.5cm}}{%
    \tableauHeader{Lab} &
    \tableauHeader{Sujet} &
    \tableauHeader{Chapitre}
}
IIN-LAB-01 & Scellés et chaîne de custody~: l'affaire MBONGO au tribunal & 8 \\
IIN-LAB-02 & Vérifier une acquisition~: image segmentée et journal FTK Imager & 9 \\
IIN-LAB-03 & Métadonnées~: qui a écrit la lettre anonyme~? & 3 \\
IIN-LAB-04 & Graphe des flux~: trouver le compte pivot & 4 \\
IIN-LAB-05 & Empreintes~: doublons et fichiers connus & 4 \\
IIN-LAB-06 & PIU~: Quality Gate~2 et hypothèses concurrentes & 6 \\
IIN-LAB-07 & Ordre de volatilité~: relire une intervention & 7 \\
IIN-LAB-08 & Coopération internationale~: conserver avant de divulguer & 10 \\
IIN-LAB-09 & Qualifier les faits~: de l'acte à l'article & 12 \\
IIN-LAB-10 & Relire un rapport~: constatations, analyse, biais & 13 \\
\end{tableaucours}
\vspace{-0.8em}
\begin{center}{\small\itshape Les laboratoires du cours dans NDAP SUKULU}\end{center}
